\par\Needspace{6\baselineskip}\medskip
\hypertarget{native-094-article072}{}\pdfbookmark[2]{72.}{article-072}
\begin{center}72.\end{center}
Finem huic Sectioni imponent quaedam problemata, quae in determinatione orbitarum planetarum et cometarum usum frequentem praestant. Ac primo quidem ad parallaxem reveniemus, a qua locum observatum liberare in art. 70 docuimus. Talis reductio ad centrum terrae, quum planetae distantiam a terra proxime saltem notam supponat, institui nequit, quoties planetae observati orbita omnino adhuc incognita est. Attamen in hoc quoque casu finem saltem eundem assequi licet, cuius caussa reductio ad centrum terrae suscipitur, ideo scilicet, quod hoc centro in plano eclipticae iacente vel iacere supposito plures formulae maiorem simplicitatem et concinnitatem nanciscuntur, quam si observatio ad punctum extra planum eclipticae referretur. Hoc itaque respectu nihil interest, utrum observatio ad centrum terrae an ad quodvis aliud punctum in plano eclipticae reducatur. Iam patet, si ad hunc finem punctum intersectionis plani eclipticae cum recta a planeta ad locum verum observationis ducta eligatur, observationem ipsam nulla prorsus reductione opus habere, quum planeta ex omnibus punctis illius rectae perinde videatur\footnote{Si ultima praecisio desideraretur, intervallum temporis, intra quod lumen a vero loco observationis ad fictum seu ab hoc ad illum delabitur, tempori proposito vel addere vel inde subducere oporteret, siquidem de locis aberratione affectis agitur: sed haec differentia vix ullius momenti esse potest, nisi latitudo perparva fuerit.}: quamobrem \SourceLeaf{095}hoc punctum quasi locum fictum observationis pro vero substituere licebit. Situm illius puncti sequenti modo determinamus.

Sit corporis coelestis longitudo $\lambda$, latitudo $\beta$, distantia $\Delta$, omnia respectu loci veri observationis in terrae superficie, cuius zenith respondeat longitudo $l$, latitudo $b$; porro sit $\pi$ semidiameter terrae, $L$ longitudo heliocentrica centri terrae, $B$ eiusdem latitudo, $R$ eiusdem distantia a Sole; denique $L'$ longitudo heliocentrica loci ficti, $R'$ ipsius distantia a Sole, $\Delta+\delta$ ipsius distantia a corpore coelesti. Tunc nullo negotio eruentur aequationes sequentes, denotante $N$ angulum arbitrarium:
\[
\begin{aligned}
R'\cos(L'-N)+\delta\cos\beta\cos(\lambda-N)
  &=R\cos B\cos(L-N)\\
  &\quad+\pi\cos b\cos(l-N)\\
R'\sin(L'-N)+\delta\cos\beta\sin(\lambda-N)
  &=R\cos B\sin(L-N)\\
  &\quad+\pi\cos b\sin(l-N)\\
\delta\sin\beta&=R\sin B+\pi\sin b.
\end{aligned}
\]
Statuendo itaque
\[
\begin{aligned}
\mathrm{I.}\qquad &(R\sin B+\pi\sin b)\operatorname{cotang}\beta=\mu,
\end{aligned}
\]
erit
\[
\begin{aligned}
\mathrm{II.}\qquad R'\cos(L'-N)&=R\cos B\cos(L-N)\\
 &\quad+\pi\cos b\cos(l-N)-\mu\cos(\lambda-N)\\
\mathrm{III.}\qquad R'\sin(L'-N)&=R\cos B\sin(L-N)\\
 &\quad+\pi\cos b\sin(l-N)-\mu\sin(\lambda-N)\\
\mathrm{IV.}\qquad \delta&=\frac{\mu}{\cos\beta}\mathpunct{\cdot}
\end{aligned}
\]

Ex aequationibus II, III determinari poterunt $R'$ et $L'$, ex IV intervallum temporis tempori observationis addendum quod erit minutis secundis $=493\delta$.

Hae aequationes sunt exactae et generales, poteruntque tunc quoque adhiberi, ubi pro plano eclipticae aequatore substituto $L$, $L'$, $l$, $\lambda$ designant ascensiones rectas, $B$, $b$, $\beta$ declinationes. Sed in casu de quo hic potissimum agimus, scilicet ubi locus fictus in ecliptica situs esse debet, exiguitas quantitatum $B$, $\pi$, $L'-L$ adhuc quandam formularum praecedentium contractionem permittit. Poterit enim pro $\pi$ assumi parallaxis media solaris, $B$ pro $\sin B$, $1$ pro $\cos B$ et $\cos(L'-L)$, $L'-L$ pro $\sin(L'-L)$. Ita faciendo $N=L$, formulae praecedentes assumunt formam sequentem:
\[
\begin{aligned}
\mathrm{I.}\qquad \mu&=(RB+\pi\sin b)\operatorname{cotang}\beta\\
\mathrm{II.}\qquad R'&=R+\pi\cos b\cos(l-L)-\mu\cos(\lambda-L)\\
\mathrm{III.}\qquad L'-L&=\frac{\pi\cos b\sin(l-L)-\mu\sin(\lambda-L)}{R'}\mathpunct{\cdot}
\end{aligned}
\]
\SourceLeaf{096}Proprie quidem hic $B$, $\pi$, $L'-L$ in partibus radii exprimendi sunt; sed patet, si illi anguli in minutis secundis exprimantur, aequationes I, III sine mutatione retineri posse, pro II autem substitui debere
\[
R'=R+\frac{\pi\cos b\cos(l-L)-\mu\cos(\lambda-L)}{206265''}\mathpunct{\cdot}
\]
Ceterum in formula III pro denominatore $R'$ absque errore sensibili semper adhibere licebit $R$. Reductio temporis autem, angulis in minutis secundis expressis, fiet
\[
=\frac{493^{\mathrm{s}}\mathpunct{.}\mu}{206265''\mathpunct{.}\cos\beta}\mathpunct{\cdot}
\]

